% ==================================================================
% id:       app-009
% topic:    שימושי הנגזרת
% subtopic: בעיות קיצון
% solution: full
% ==================================================================
מצאו את המרחק הקצר ביותר בין הנקודה \(M(p, p)\) ונקודות הפרבולה \(y^2 = 2px\).

\hint{תארו את נקודות הפרבולה בעזרת $y$ בלבד, ומזערו את ריבוע המרחק במקום את המרחק עצמו.}
\begin{solbox}
תהא
$(x_0,y_0)$
נקודה על הפרבולה, ניתן לרשום
$$.(x_0,y_0)=\left(\frac{y_0^2}{2p},y_0\right)$$
כזכור משיעורי גיאומטריה בתיכון, המרחק בין שתי נקודות
$(a,b)$
ו-
$(c,d)$
מחושב ע"י
$$.\sqrt{(a-c)^2+(b-d)^2}$$
מכאן המרחק בין הנקודה
$M$
לנקודה
$(x_0,y_0)$
נתון ע"י
$$.\sqrt{\left(\frac{y_0^2}{2p}-p\right)^2+(y_0-p)^2}$$
עלינו למזער את הפונקציה
$$.g(x)=\sqrt{\left(\frac{x^2}{2p}-p\right)^2+(x-p)^2},\ x\in\mathbb{R}$$
נשים לב כי
$x_1$
מספקת מינימום מוחלט ל-
$g(x)$
אמ"מ
$x_1$
מספקת מינימום מוחלט ל-
$g^2(x)$
(האם תוכלו להראות זאת?).
לאור הנאמר, נעבוד עם פונקציה
$$,f(x)=\left(\frac{x^2}{2p}-p\right)^2+(x-p)^2,\ x\in\mathbb{R}$$
לאחר פתיחת סוגריים נקבל
$$.f(x)=\frac{x^4}{4p^2}-2px+2p^2$$
נגזור ונשווה ל-
$0$:
$$,0=f'(x)=\frac{x^3}{p^2}-2p$$
נסיק כי
$a=\sqrt[3]{2}p$.
מתקיים
$$.f(a)=p^2\cdot\frac{4-3\sqrt[3]{2}}{2}$$
נסמן
$a=p^2\cdot\frac{4-3\sqrt[3]{2}}{2}$.
נראה שזוהי נקודת מינימום מוחלט של
$f$.
אכן, נניח בשלילה כי
$x_0\in\mathbb{R}$
היא מינימום מוחלט.
מכיוון ש-
$$,\lim_{x\to-\infty}f(x)=\infty,\ \lim_{x\to\infty}f(x)=\infty$$
קיימים
$M_1<0$
ו-
$M_2>0$
כך ש-
$f(x)>f(a)$
לכל
$x<M_1$
ולכל
$x>M_2$.
נתבונן בקטע סגור
$[M_1-1,M_2+1]$,
 מכיוון ש-
$f$
רציפה וגזירה ב-
$[M_1,M+1]$
היא משיגה את המינימום המוחלט שלה או בקצוות הקטע או בנקודה
שמאפסת את נגזרתה.
לפי ההנחה שלנו
$$f(M_1-1)>f(a)>f(x_0)$$
ו-
$$,f(M_2+1)>f(a)>f(x_0)$$
לכן נסיק כי
$f'(x_0)=0$,
אבל זה לא מתקיים, לכן
$x_0\neq a$
לא יכולה להיות נקודת שמקיימת
$f(x_0)<f(a)$,
מה שמראה כי
$a$
היא אכן נקודת מינימום.
מכאן המרחק המינימלי הוא
$$.\sqrt{f(a)}=|p|\cdot\sqrt{\frac{4-3\sqrt[3]{2}}{2}}$$
\end{solbox}
